\section{Rational Man and Moral Law}

\label{sec:RationalMalMoralLaw}

\begin{quotex}
Today the meaning of the moral as cosmic worth has been lost, and up to now, it has been reduced to nothing more than the corroboration and canonisation of deficiency, weakness, and fear — of the plague of beautiful sentiments, of noble virtues, of holy ideals — that is, of the fundamentally immoral. \flright{\textsc{Julius Evola}}

\end{quotex}
\paragraph{St Anthony the Great}
St Anthony has already been dealt with in Section~\ref{sec:St Anthony and Practical Reason}, so we will focus just on his concept of the “rational man”. The rational man is not just a learned man or scholar, but rather someone who not only knows good from bad, but also has the power to act on it. The rational faculty needs to be developed. In appearance, Anthony seems to be recommending conventional morality when he talks about overcoming irrational passions and desires or orienting oneself toward the eternal and divine rather than the temporal and material. Yet these are more than moral strictures, they are rational, properly so called.

The rational faculty unites man with divine power, power over oneself, one's thoughts, one's acts. As such, it is a faculty that must be developed. For Anthony, the irrational man is in a state of darkness and devoid of divine light. That is, the rational faculty is only virtual and not actual. Since man by definition is rational, such men are not, strictly speaking, even men. This brings Anthony close to the Gnostic idea of the three classes of men [Hylic, Psychic, and Pneumatic], and only the actualized rational man is pneumatic. 

There is a potential consequence to this point of view that does not seem to have been realized by Anthony's commentators, though it has by Nietzsche. Eric Voegelin writes, in his discussion of what Reason (or rationality) is:

\begin{quotex}
The constituent of society is the Homonoia or “like-mindedness” of Everyman in a community formed through recognition of the reason common to all men. In Aristotle, if love within the community is not based upon regard for the divinity of reason in the other man, then the political friendship on which a well-ordered community depends cannot exist. The source of the Christian notion of “human dignity” is the common divinity in all men. Nietzsche perceived that if that is surrendered then there is no reason to love anybody, one consequence of which is the loss of the sense and force of obligation in society and, hence, of its cohesiveness. 

\end{quotex}
But according to Anthony, only the rational man is in union with God (or Nous in Platonic or Aristotelian terms), and reason is not at all common to all men. Obviously, the modernist project of deriving human dignity from material and biological considerations has failed (it is unreasonable), with the resulting loss of cohesiveness in Western societies. So, if biologism cannot justify egalitarianism, then neither can Christianity, since full human dignity may only be virtual. Another way of putting it is to say that although in essence there may be a common divinity in men, it is not necessarily so in existence.

Yet, Anthony insists on comparing the irrational man to brute animals, whether or not the rational faculty is merely undeveloped or, perhaps, does not even exist at all. The full implications for ethics and political science have not yet been drawn out.

\paragraph{Julius Evola}
If the world is Will and Idea, then knowledge of ideas is insufficient for its understanding and the hope for a science of everything is futile. The act of Will brings Ideas into reality, that is, it causes a change in Being. Thus the world does not present itself as a scientific problem, but rather as a moral problem, so fundamentally, morality necessarily implies the power to act. Evola writes:

\begin{quotex}
The ethical problem and the metaphysical problem coincide: the measure of reality, as of the good, of the certain and of the true, is the perfection of actuality or power. Per virtutem et potentiam idem intelligo. [Spinoza, Ethics, “by virtue and power I mean the same thing.”] The good or virtue is potency, evil impotence. There is no other evil except insufficiency and weakness, no other good except the will that is unconditionally sufficient in itself, and therefore free to be and do what it wills.

\end{quotex}
Conventional morality is of two types. One is the idea of a war against the body or parts of the soul. The other is intellectual and involves finding a moral system that demands adherence. Neither solves the problem of power or change of being.

More recent attempts — such as that of Ken Wilber — have postulated a scale of moralities, as one “evolves” from stage to stage. Such attempts are clearly self-serving and unprovable, not to mention incompatible with a true moral will. Evola anticipated such scales when he wrote:

\begin{quotex}
The life of the spirit is essentially that of freedom. It makes no sense to speak of a transcendent law — whether natural, rational, or moral — which would determine how the stages would proceed in accordance with some incontrovertible sequence: the passage from one stage to the other is instead unconditioned, it happens if it happens, its direction is not deducible from one stage to another; it is an absolute achievement, an achievement, therefore, for which one cannot, but even must not, ask for a further reason.

\end{quotex}

The material on Anthony is taken from the booklet “Rational Man” by Constantine Cavarnos and is available from the Institute for Byzantine and Modern Greek Studies\footnote{\url{http://www.ibmgs.org/}}.

\flrightit{Posted on 2010-12-02 by Cologero}

\begin{center}* * *\end{center}

\begin{footnotesize}\begin{sffamily}
\texttt{Matt on 2010-12-02 at 12:21 said: }

Would it be valid to say that Anthony's notion of the divine rational faculty being present in man in essence but not necessarily in existence is in the same spirit as Evola's (and in the wider context of Tradition) view that the Self is present in essence but not existence?

Also, after getting more acquainted with St. Anthony's view of what makes one human and its relation to morality, I now further understand your comment modern day Christians being hesitant to include his writings in current translations of the Philokalia.

\hfill

\texttt{GF on 2010-12-09 at 00:08 said: }

So the true moral project is the struggle to make oneself exist, or in religious language, to heal the soul. Can we perhaps see ‘civic virtue’ as a corresponding exoterism?


\hfill

\texttt{Cologero on 2010-12-10 at 07:39 said: }

Yes, that was the point of the juxtaposition, to demonstrate the compatibility of Anthony with Tradition. We always need to avoid being distracted by the theology and dig down to the underlying metaphysical, or esoteric, core. I've been told by someone who has been involved with esoterism and mysteries for over 50 years that the Coptic religion (Anthony was Egyptian) is a Christian veneer over the ancient pharaonic religion. I have no way to verify that, as interesting as it sounds. (The Coptic language is based on ancient Egyptian.)


\hfill

\texttt{Cologero on 2010-12-10 at 07:43 said: }

Exactly, if “heal” means “make whole”. Not ‘civic virtue’ if you mean an external code of conduct that “I” must follow. But yes, if on the lines of synarchy, such as described in The City of the Sun.


\hfill

\texttt{Matt on 2010-12-10 at 10:49 said: }

Cologero,

Yes that is really interesting. If thats the case, Coptic initiates achieved what the Templars and Hermeticists ultimately failed to do with the catholic church.


\hfill
\end{sffamily}\end{footnotesize}
